% !TEX root = ../tcc.tex
\section{Metodologia}\label{sec:metodologia}

\subsection{Pipeline}\label{sec:pipeline}
% Documento-fonte -> Extrator -> Analista -> Redator -> Resposta final.
% Modelo, prompts, documentos e lógica dos estágios constantes entre condições.

\subsection{Unidade de análise: o claim}\label{sec:claim}
% claim_id, derived_from e a cadeia de proveniência (ex.: C31 -> C17 -> C08).

\subsection{Condições de contrato de handoff}\label{sec:condicoes}
% A (texto livre), B (tipado), C (tipado + lineage), D (tipado + lineage + evidência).
% Tabela comparativa e o controle de validade interna: em C, source_id/page são
% só metadados; apenas D recebe o trecho de evidência.

\subsection{Injeção de erros}\label{sec:injecao}
% Três classes (valor, percentual, metadado), perturbações determinísticas
% registradas como ground truth; mesmo estado corrompido serializado em A/B/C/D.

\subsection{Classificação do destino do erro}\label{sec:destino}
% Preservado / transformado / amplificado / corrigido / contido.
